\section{Theoretical Analysis of GQA Amplification}
\label{sec:theory}

In this section we formalize the empirical observation from
Section~\ref{sec:intro}: vector quantization of the KV cache degrades faster
than per-channel scalar quantization as the grouped-query-attention ratio
increases. We derive a first-order upper bound on the attention output error
under KV compression, decompose it into a per-head error variance term and a
GQA aggregation term, and show that the aggregation term scales linearly in
the group size~$g$ for VQ-style errors but only logarithmically (and at most
sublinearly) for per-channel scalar errors. The bound motivates the practical
recipe we adopt: use vector quantization at low GQA ratios, and switch to a
hybrid scheme (outlier protection or sensitivity-adaptive allocation) at high
ratios where the amplification term becomes the dominant source of error.

\subsection{Setup}
\label{sec:theory-setup}

Consider a single attention layer with $H$ query heads and $H_{kv}$ key-value
heads, where $g = H / H_{kv}$ is the GQA group size. Each KV head $j \in
\{1, \ldots, H_{kv}\}$ is shared by exactly $g$ query heads. Let
$q_h \in \mathbb{R}^{d_k}$ denote a single query vector (any position, any
head), and let $K_j \in \mathbb{R}^{N \times d_k}$, $V_j \in \mathbb{R}^{N
\times d_v}$ be the keys and values for KV head $j$ over a context of length
$N$. In practice $d_k = d_v = d$.

The uncompressed attention output for query head $h$ (assume it belongs to
group $j$) is
\begin{equation}
\label{eq:attn-clean}
o_h = \mathrm{softmax}\!\left( \tfrac{1}{\sqrt{d}} q_h K_j^\top \right) V_j.
\end{equation}
Under KV compression, the stored $K_j, V_j$ are replaced by their
reconstructions $\tilde K_j = K_j + \Delta^K_j$ and
$\tilde V_j = V_j + \Delta^V_j$, where $\Delta^K_j$ and $\Delta^V_j$ are the
compression residuals.  The noisy attention output is
\begin{equation}
\label{eq:attn-noisy}
\tilde o_h = \mathrm{softmax}\!\left( \tfrac{1}{\sqrt{d}} q_h \tilde K_j^\top \right) \tilde V_j.
\end{equation}
We are interested in bounding the per-layer output perturbation
$\|\tilde o_h - o_h\|_2$ and, in aggregate, the full layer output error
summed over all $H$ query heads.

\subsection{Per-Query Perturbation Bound}
\label{sec:theory-perquery}

Define the attention logit vector $\ell_h = q_h K_j^\top / \sqrt{d}
\in \mathbb{R}^N$ and its perturbed counterpart $\tilde \ell_h =
q_h \tilde K_j^\top / \sqrt{d}$. The logit perturbation at position $t$ is
\begin{equation}
\label{eq:delta-logit}
\delta_{h,t} := \tilde \ell_{h,t} - \ell_{h,t} = \frac{1}{\sqrt{d}}\, q_h^\top \Delta^K_{j,t},
\end{equation}
so $|\delta_{h,t}| \le \|q_h\|_2 \|\Delta^K_{j,t}\|_2 / \sqrt{d}$. Under the
standard assumption $\|q_h\|_2 = \Theta(\sqrt{d})$ (queries are near the unit
sphere after RoPE scaled by $\sqrt{d}$, consistent with the preactivation
distribution of a trained transformer), we have
\begin{equation}
\label{eq:bound-logit}
|\delta_{h,t}| \le c_q \cdot \|\Delta^K_{j,t}\|_2,
\end{equation}
for a model-dependent constant $c_q \approx 1$.

The softmax is 1-Lipschitz in its input with respect to the total-variation
norm on the output distribution \cite{gao2017properties}:
\[
\|\mathrm{softmax}(\tilde \ell_h) - \mathrm{softmax}(\ell_h)\|_1
\le \|\tilde \ell_h - \ell_h\|_\infty.
\]
Combining this with~\eqref{eq:bound-logit} and the triangle inequality applied
to the value perturbation gives a first-order bound on the output:
\begin{equation}
\label{eq:bound-output}
\|\tilde o_h - o_h\|_2
\le \underbrace{\|V_j\|_{\mathrm{op}} \cdot c_q \cdot \max_t \|\Delta^K_{j,t}\|_2}_{\text{key-side term}}
\; + \;
\underbrace{\max_t \|\Delta^V_{j,t}\|_2}_{\text{value-side term}}
\; + \; O(\varepsilon^2),
\end{equation}
where $\|V_j\|_{\mathrm{op}}$ is the operator norm of $V_j$ and
$\varepsilon = \max(\|\Delta^K\|_\infty, \|\Delta^V\|_\infty)$.
The first-order approximation is valid in the small-noise regime that
compressed caches typically operate in ($\varepsilon \ll 1$ in unit-norm KV
space), which we verify empirically in Section~\ref{sec:qwen-rescue}: even
when the compressed PPL is large (42 on Qwen2.5-3B), the per-vector MSE is
still $\sim 10^{-2}$, well below unity.

Equation~\eqref{eq:bound-output} shows that the per-query perturbation scales
linearly with the reconstruction error magnitude on a \emph{single} KV head.
This bound does not yet capture the GQA effect.

\subsection{GQA Aggregation}
\label{sec:theory-gqa}

A single compressed KV head $j$ is consumed by all $g$ query heads in its
group. The \emph{layer-level} output error in the residual stream is the sum
of the per-head errors:
\begin{equation}
\label{eq:layer-error}
\mathcal{E}_{\mathrm{layer}}^2 = \mathbb{E}\!\left[ \Big\| \sum_{j=1}^{H_{kv}} \sum_{h \in \mathrm{group}(j)} (\tilde o_h - o_h) \Big\|_2^2 \right].
\end{equation}
The nature of the inner sum over $h \in \mathrm{group}(j)$ depends on the
\emph{structure} of the compression error $\Delta^K_j, \Delta^V_j$.

\paragraph{Case 1 — Vector-quantization (VQ) errors.}
For TurboQuant-style compression, the error $\Delta^K_j$ at a given position
$t$ is a vector-valued codebook quantization residual, fully correlated
across the $d$ channels and identical for every query head that reads from
$K_j$. Writing $\delta_h^{(j)} = \tilde o_h - o_h$, we have
\[
\sum_{h \in \mathrm{group}(j)} \delta_h^{(j)} = g \cdot \bar\delta^{(j)} + \text{(cross-head noise)},
\]
where $\bar\delta^{(j)}$ is the mean error across the $g$ query heads in the
group, $O(1)$ in magnitude by~\eqref{eq:bound-output}. The cross-head noise
averages out in expectation for random queries, but the coherent term grows
linearly in $g$. Summing $\mathcal{E}_{\mathrm{layer}}^2$ over groups gives
\begin{equation}
\label{eq:vq-scaling}
\mathcal{E}_{\mathrm{layer}}^{\mathrm{VQ}} = O\!\left( g \cdot \sqrt{H_{kv}} \cdot \sigma_{\mathrm{vec}} \right),
\end{equation}
where $\sigma_{\mathrm{vec}}^2 = \mathbb{E}\|\Delta^K_j\|_2^2$ is the
per-vector MSE of the vector quantizer. The layer error grows
\emph{linearly} in the GQA ratio $g$.

\paragraph{Case 2 — Per-channel scalar errors (KIVI-style).}
For KIVI, each element $\Delta^K_{j,t,c}$ at channel $c$ is an independent
scalar quantization residual with per-channel variance
$\sigma_{\mathrm{ch}}^2$. The per-vector norm is
$\|\Delta^K_{j,t}\|_2^2 \approx d \cdot \sigma_{\mathrm{ch}}^2$ in expectation,
but the contributions to the attention logit from individual channels are
\emph{uncorrelated with each other and with the query pattern}. The
per-query error then behaves like a sum of $d$ random contributions
(central-limit heuristic), and the cross-head sum behaves like a sum of $g$
sub-Gaussian random variables whose variance grows only $\sqrt{g}$:
\begin{equation}
\label{eq:kivi-scaling}
\mathcal{E}_{\mathrm{layer}}^{\mathrm{KIVI}} = O\!\left( \sqrt{g} \cdot \sqrt{H_{kv}} \cdot \sqrt{d} \cdot \sigma_{\mathrm{ch}} \right).
\end{equation}
The layer error grows as $\sqrt{g}$ in the GQA ratio, not linearly.

\subsection{Crossover Inequality}
\label{sec:theory-crossover}

Comparing~\eqref{eq:vq-scaling} and~\eqref{eq:kivi-scaling}, VQ yields a
lower layer-output error than KIVI when
\begin{equation}
\label{eq:crossover}
\sigma_{\mathrm{vec}}^2 \;\lesssim\; \frac{d}{g} \cdot \sigma_{\mathrm{ch}}^2.
\end{equation}
At fixed bit-width and a reasonably well-behaved KV distribution, TurboQuant's
rotation-then-Lloyd-Max pipeline achieves $\sigma_{\mathrm{vec}}^2 \approx
\beta_b \cdot \sigma_{\mathrm{ch}}^2$ for some constant $\beta_b$ that depends
on the bit-width $b$. The crossover condition becomes
\begin{equation}
\label{eq:crossover-b}
\beta_b \;\lesssim\; d / g,
\end{equation}
which is easily satisfied at $b = 4$, $d = 128$, $g = 4$ (tolerance factor
$\approx 32$), but becomes tight at $g = 8$ (tolerance $\approx 16$) and
can be violated on models whose KV distributions concentrate mass in a few
outlier channels, inflating $\beta_b$.

This matches our empirical observations:

\begin{itemize}
  \item \textbf{Gemma-3-27B (GQA $2{:}1$)}: crossover tolerance $\approx 64$.
  TurboQuant K4V4 is essentially lossless ($+0.7\%$ PPL) while KIVI-4
  degrades by $+22.4\%$.
  \item \textbf{Mistral-7B, Llama-3.1-8B (GQA $4{:}1$)}: crossover tolerance
  $\approx 32$. TurboQuant still wins by roughly an order of magnitude.
  \item \textbf{Qwen2.5-3B (GQA $8{:}1$)}: crossover tolerance $\approx 16$.
  Combined with Qwen's outlier-heavy layer-0 KV distribution (which inflates
  $\beta_b$), the inequality is violated and TurboQuant collapses. Outlier
  protection (Section~\ref{sec:outlier}) directly attacks $\beta_b$ by
  isolating the high-magnitude channels that drive it; adaptive allocation
  (Section~\ref{sec:adaptive}) attacks the same problem from the other side
  by raising the effective $b$ on the most sensitive layers.
\end{itemize}

\subsection{Limitations of the Analysis}
\label{sec:theory-limitations}

Our bound is a first-order worst-case perturbation analysis and omits several
effects that matter in practice. First, the softmax is not globally
Lipschitz in $\ell_\infty$; the bound tightens near uniform distributions and
loosens when attention is strongly peaked. Second, the residual stream
aggregates errors across layers, and deep layers can amplify or cancel errors
nonlinearly through LayerNorm and MLP blocks, a higher-order effect we do not
attempt to model. Third, we treat $\sigma_{\mathrm{vec}}$ and
$\sigma_{\mathrm{ch}}$ as fixed; in reality both depend on the KV
distribution and fluctuate across layers (which is precisely why adaptive
allocation helps). Nevertheless, the linear-in-$g$ dependence of the VQ
error is a structural property of shared-group codebook quantization and is
independent of these refinements. Any tighter bound we could derive would
still exhibit the same asymptotic scaling, and so the qualitative conclusion
--- VQ is a better fit for low-$g$ models than high-$g$ models --- holds
robustly.
